\chapter{2006年全国I卷（文）}
\section{单选题}
\begin{problem}
已知向量 \(\bs{a}\)，\(\bs{b}\) 满足 \(|\bs{a}| = 1\)，\(|\bs{b}| = 4\)，且 \(\bs{a} \cdot \bs{b} = 2\)，则 \(\bs{a}\) 与 \(\bs{b}\) 夹角为（\quad）

\choices
    {\(\frac{\pi}{6}\)}
    {\(\frac{\pi}{4}\)}
    {\(\frac{\pi}{3}\)}
    {\(\frac{\pi}{2}\)}
\end{problem}

\begin{answer}
C
\end{answer}

\begin{solution}
设 \(\bs{a}\) 与 \(\bs{b}\) 的夹角为 \(\theta\). 由数量积公式，
\(\bs{a}\cdot\bs{b}=|\bs{a}|\,|\bs{b}|\cos\theta\)，所以 \(2=1\times4\cos\theta\)，得 \(\cos\theta=\frac{1}{2}\). 因为 \(0\leqslant\theta\leqslant\pi\)，故 \(\theta=\frac{\pi}{3}\)，所以选 C.
\end{solution}

\begin{problem}
设集合 \(M = \{x \mid x^2 - x < 0\}\)，\(N = \{x \mid |x| < 2\}\)，则（\quad）

\choices
    {\(M \cap N = \varnothing\)}
    {\(M \cap N = M\)}
    {\(M \cup N = M\)}
    {\(M \cup N = \R\)}
\end{problem}

\begin{answer}
B
\end{answer}

\begin{solution}
由 \(x^{2}-x<0\) 得 \(0<x<1\)，所以 \(M=(0,1)\). 又由 \(|x|<2\) 得 \(N=(-2,2)\). 因而 \(M\subseteq N\)，从而 \(M\cap N=M\)，所以选 B.
\end{solution}

\begin{problem}
已知函数 \( y = \e^{x} \) 的图象与函数 \( y = f(x) \) 的图象关于直线 \( y = x \) 对称，则（\quad）

\choices
    {\(f(2x) = \e^{2x}\)（\(x\in \R\)）}
    {\(f(2x) = \ln 2\cdot \ln x\)（\(x > 0\)）}
    {\(f(2x) = 2\e^{x}\)（\(x\in \R\)）}
    {\(f(2x) = \ln x + \ln 2\)（\(x > 0\)）}
\end{problem}

\begin{answer}
D
\end{answer}

\begin{solution}
关于直线 \(y=x\) 对称的两个函数互为反函数. \(y=\e^{x}\) 的反函数为 \(y=\ln x\)，定义域为 \(x>0\). 因此 \(f(2x)=\ln(2x)=\ln x+\ln2\)，且 \(x>0\)，所以选 D.
\end{solution}

\begin{problem}
双曲线 \(mx^{2} + y^{2} = 1\) 的虚轴长是实轴长的 \(2\) 倍，则 \(m =\)（\quad）

\choices
    {\(-\frac{1}{4}\)}
    {\(-4\)}
    {\(4\)}
    {\(\frac{1}{4}\)}
\end{problem}

\begin{answer}
A
\end{answer}

\begin{solution}
双曲线方程可写为 \(y^{2}-(-m)x^{2}=1\)，故必须有 \(m<0\)，其实轴半长为 \(1\)，虚轴半长为 \(\frac{1}{\sqrt{-m}}\). 题设给出虚轴长是实轴长的 \(2\) 倍，即虚轴半长是实轴半长的 \(2\) 倍，所以 \(\frac{1}{\sqrt{-m}}=2\). 解得 \(m=-\frac{1}{4}\)，所以选 A.
\end{solution}

\begin{problem}
设 \(S_{n}\) 是等差数列 \(\{a_{n}\}\) 的前 \(n\) 项和，若 \(S_{7} = 35\)，则 \(a_{4} =\)（\quad）

\choices
    {\(8\)}
    {\(7\)}
    {\(6\)}
    {\(5\)}
\end{problem}

\begin{answer}
D
\end{answer}

\begin{solution}
等差数列的前 \(7\) 项关于第 \(4\) 项成对对称，因此 \(S_{7}=7a_{4}\). 由 \(S_{7}=35\) 得 \(a_{4}=5\)，所以选 D.
\end{solution}

\begin{problem}
函数 \( f(x) = \tan \left(x + \frac{\pi}{4}\right) \) 的单调增区间为（\quad）

\choices
    {\(\left(k\pi -\frac{\pi}{2},k\pi +\frac{\pi}{2}\right),k\in \Z\)}
    {\((k\pi ,(k + 1)\pi)\)，\(k\in \Z\)}
    {\(\left(k\pi -\frac{3\pi}{4},k\pi +\frac{\pi}{4}\right)\)，\(k\in \Z\)}
    {\(\left(k\pi -\frac{\pi}{4},k\pi +\frac{3\pi}{4}\right)\)，\(k\in \Z\)}
\end{problem}

\begin{answer}
C
\end{answer}

\begin{solution}
正切函数在每个区间 \(\left(-\frac{\pi}{2}+k\pi,\frac{\pi}{2}+k\pi\right)\) 上单调递增. 对 \(f(x)=\tan\left(x+\frac{\pi}{4}\right)\)，应满足
\(x+\frac{\pi}{4}\in\left(-\frac{\pi}{2}+k\pi,\frac{\pi}{2}+k\pi\right)\)，即 \(x\in\left(k\pi-\frac{3\pi}{4},k\pi+\frac{\pi}{4}\right)\)，其中 \(k\in\Z\). 所以选 C.
\end{solution}

\begin{problem}
从圆 \(x^{2} - 2x + y^{2} - 2y + 1 = 0\) 外一点 \(P(3,2)\) 向这个圆作两条切线，则两切线夹角的余弦值为（\quad）

\choices
    {\(\frac{1}{2}\)}
    {\(\frac{3}{5}\)}
    {\(\frac{\sqrt{3}}{2}\)}
    {\(0\)}
\end{problem}

\begin{answer}
B
\end{answer}

\begin{solution}
圆的方程化为 \(\left(x-1\right)^{2}+\left(y-1\right)^{2}=1\)，圆心为 \(O(1,1)\)，半径为 \(1\). 由 \(P(3,2)\) 得 \(OP=\sqrt{5}\). 设两切点为 \(U\)，\(V\)，则 \(OU\perp PU\)，在直角三角形 \(POU\) 中，\(PU=\sqrt{OP^{2}-OU^{2}}=2\). 设 \(\angle UPO=\alpha\)，则 \(\cos\alpha=\frac{PU}{PO}=\frac{2}{\sqrt{5}}\). 两条切线关于 \(PO\) 对称，所以 \(\angle UPV=2\alpha\)，从而 \(\cos\angle UPV=2\cos^{2}\alpha-1=2\times\frac{4}{5}-1=\frac{3}{5}\)，所以选 B.
\end{solution}

\begin{problem}
\(\triangle ABC\) 的内角 \(A\)，\(B\)，\(C\) 的对边分别为 \(a\)，\(b\)，\(c\)，若 \(a\)，\(b\)，\(c\) 成等比数列，且 \(c = 2a\)，则 \(\cos B =\)（\quad）

\choices
    {\(\frac{1}{4}\)}
    {\(\frac{3}{4}\)}
    {\(\frac{\sqrt{2}}{4}\)}
    {\(\frac{\sqrt{2}}{3}\)}
\end{problem}

\begin{answer}
B
\end{answer}

\begin{solution}
设等比数列 \(a\)，\(b\)，\(c\) 的公比为 \(q\). 因为三边长为正数且 \(c=2a\)，有 \(q^{2}=\frac{c}{a}=2\)，故 \(b=\sqrt{2}a\). 由余弦定理，
\(\cos B=\frac{a^{2}+c^{2}-b^{2}}{2ac}=\frac{a^{2}+4a^{2}-2a^{2}}{4a^{2}}=\frac{3}{4}\)，所以选 B.
\end{solution}

\begin{problem}
已知各顶点都在一个球面上的正四棱柱高为 \(4\)，体积为 \(16\)，则这个球的表面积是（\quad）

\choices
    {\(16\pi\)}
    {\(20\pi\)}
    {\(24\pi\)}
    {\(32\pi\)}
\end{problem}

\begin{answer}
C
\end{answer}

\begin{solution}
设正四棱柱底面边长为 \(s\). 由体积条件，\(s^{2}\times4=16\)，得 \(s=2\). 该棱柱的空间对角线长为 \(\sqrt{2^{2}+2^{2}+4^{2}}=2\sqrt{6}\)，外接球半径为 \(\sqrt{6}\). 因而球的表面积为 \(4\pi\left(\sqrt{6}\right)^{2}=24\pi\)，所以选 C.
\end{solution}

\begin{problem}
在 \(\left(x - \frac{1}{2x}\right)^{10}\) 的展开式中，\(x^4\) 的系数为（\quad）

\choices
    {\(-120\)}
    {\(120\)}
    {\(-15\)}
    {\(15\)}
\end{problem}

\begin{answer}
C
\end{answer}

\begin{solution}
展开式的通项为 \(\mathrm C_{10}^{k}x^{10-k}\left(-\frac{1}{2x}\right)^{k}=\mathrm C_{10}^{k}\frac{(-1)^{k}}{2^{k}}x^{10-2k}\). 要得到 \(x^{4}\)，需有 \(10-2k=4\)，所以 \(k=3\). 所求系数为 \(\mathrm C_{10}^{3}\frac{(-1)^{3}}{2^{3}}=-\frac{120}{8}=-15\)，所以选 C.
\end{solution}

\begin{problem}
抛物线 \( y = -x^{2} \) 上的点到直线 \( 4x + 3y - 8 = 0 \) 距离的最小值是（\quad）

\choices
    {\(\frac{4}{3}\)}
    {\(\frac{7}{5}\)}
    {\(\frac{8}{5}\)}
    {\(3\)}
\end{problem}

\begin{answer}
A
\end{answer}

\begin{solution}
取抛物线上的点 \(\left(x,-x^{2}\right)\). 它到直线 \(4x+3y-8=0\) 的距离为
\(\frac{|4x-3x^{2}-8|}{5}=\frac{3x^{2}-4x+8}{5}\)，因为 \(3x^{2}-4x+8>0\). 配方得
\(\frac{3x^{2}-4x+8}{5}=\frac{3\left(x-\frac{2}{3}\right)^{2}+\frac{20}{3}}{5}\geqslant\frac{4}{3}\). 当 \(x=\frac{2}{3}\) 时等号成立，所以最小值为 \(\frac{4}{3}\)，选 A.
\end{solution}

\begin{problem}
用长度分别为 \(2\)，\(3\)，\(4\)，\(5\)，\(6\)（单位：\(\mathrm{cm}\)）的 \(5\) 根细木棒围成一个三角形（允许连接，但不允许折断），能够得到的三角形的最大面积为（\quad）

\choices
    {\(8\sqrt{5}\mathrm{~cm}^{2}\)}
    {\(6\sqrt{10}\mathrm{~cm}^{2}\)}
    {\(3\sqrt{55}\mathrm{~cm}^{2}\)}
    {\(20\mathrm{~cm}^{2}\)}
\end{problem}

\begin{answer}
B
\end{answer}

\begin{solution}
把木棒连接后，每一条三角形边由若干根木棒首尾相接而成，因此只需把总长 \(20\) 分成三个满足三角形不等式的边长. 分组大小只能是 \(3+1+1\) 或 \(2+2+1\).

当分组为 \(3+1+1\) 时，满足三角形不等式的边长只有 \(5\)，\(6\)，\(9\)，其面积为 \(10\sqrt{2}\). 当分组为 \(2+2+1\) 时，去掉单根木棒后配成两组，满足三角形不等式的边长依次为 \((2,9,9)\)、\((3,8,9)\)、\((4,7,9)\)、\((4,8,8)\)、\((5,6,9)\)、\((5,7,8)\)、\((6,6,8)\)、\((6,7,7)\). 这些三角形的半周长均为 \(10\)，由海伦公式得到面积依次为 \(4\sqrt{5}\)、\(2\sqrt{35}\)、\(6\sqrt{5}\)、\(4\sqrt{15}\)、\(10\sqrt{2}\)、\(10\sqrt{3}\)、\(8\sqrt{5}\)、\(6\sqrt{10}\). 比较这些面积的平方可知最大的是 \(6\sqrt{10}\). 例如该分组由长度为 \(6\) 的木棒、长度为 \(2\) 和 \(5\) 的两根木棒以及长度为 \(3\) 和 \(4\) 的两根木棒组成，确实可行，所以选 B.
\end{solution}

\section{填空题}
\begin{problem}
已知函数 \( f(x) = a - \frac{1}{2^x + 1} \)，若 \( f(x) \) 为奇函数，则 \( a = \)\fillinblank{}.
\end{problem}

\begin{answer}
\(\frac{1}{2}\)
\end{answer}

\begin{solution}
奇函数满足 \(f(0)=0\)，所以 \(a-\frac{1}{2^{0}+1}=0\)，即 \(a=\frac{1}{2}\). 反之，取 \(a=\frac{1}{2}\) 时，令 \(t=2^{x}>0\)，则
\(f(x)=\frac{1}{2}-\frac{1}{t+1}=\frac{t-1}{2(t+1)}\)，而 \(2^{-x}=\frac{1}{t}\)，故 \(f(-x)=\frac{1-t}{2(1+t)}=-f(x)\). 因而该函数确为奇函数，填 \(\frac{1}{2}\).
\end{solution}

\begin{problem}
已知正四棱锥的体积为 \(12\)，底面对角线长为 \(2\sqrt{6}\)，则侧面与底面所成的二面角等于\fillinblank{}.
\end{problem}

\begin{answer}
\(\frac{\pi}{3}\)
\end{answer}

\begin{solution}
设正四棱锥底面中心为 \(O\)，底面一边的中点为 \(M\)，顶点为 \(P\). 由底面对角线长为 \(2\sqrt{6}\)，得底面边长为 \(2\sqrt{3}\)，底面积为 \(12\). 由体积 \(\frac{1}{3}\times12\times PO=12\)，得 \(PO=3\). 又 \(OM\) 是底面边长的一半，故 \(OM=\sqrt{3}\). 在截面直角三角形 \(POM\) 中，侧面与底面的二面角为 \(\angle PMO\)，且 \(\tan\angle PMO=\frac{PO}{OM}=\sqrt{3}\). 因为二面角为锐角，故其大小为 \(\frac{\pi}{3}\).
\end{solution}

\begin{problem}
设 \(z = 2y - x\)，式中变量 \(x\)，\(y\) 满足下列条件：\(\begin{cases}
2x - y \geqslant -1,\\
3x + 2y \leqslant 23,\\
y \geqslant 1,
\end{cases}\)，则 \(z\) 的最大值为\fillinblank{}.
\end{problem}

\begin{answer}
11
\end{answer}

\begin{solution}
约束条件等价于 \(y\leqslant2x+1\)，\(y\leqslant\frac{23-3x}{2}\)，且 \(y\geqslant1\). 三条边界线的交点分别为 \((0,1)\)、\((7,1)\) 和 \((3,7)\). 目标函数 \(z=2y-x\) 在这个三角形区域上取最大值时只需比较三个顶点：
\(z(0,1)=2\)，\(z(7,1)=-5\)，\(z(3,7)=11\). 所以最大值为 \(11\)，填 \(11\).
\end{solution}

\begin{problem}
安排 \(7\) 位工作人员在 \(5\) 月 \(1\) 日至 \(5\) 月 \(7\) 日值班，每人值班一天，其中甲，乙二人都不安排在 \(5\) 月 \(1\) 日和 \(2\) 日，不同的安排方法共有\fillinblank{}\(\text{种}\). （用数字作答）
\end{problem}

\begin{answer}
2400
\end{answer}

\begin{solution}
甲、乙均不能安排在前两天，所以甲先有 \(5\) 种选择，乙再有 \(4\) 种选择. 甲、乙安排后，剩下的 \(5\) 位工作人员任意安排到剩余的 \(5\) 天，有 \(5!\) 种方法. 因而不同安排方法共有 \(5\times4\times5!=2400\) 种，填 \(2400\).
\end{solution}

\section{解答题}
\begin{problem}
已知 \(\{a_{n}\}\) 为等比数列，\(a_{3} = 2\)，\(a_{2} + a_{4} = \frac{20}{3}\)，求 \(\{a_{n}\}\) 的通项公式.
\end{problem}

\begin{solution}
设等比数列的公比为 \(q\). 因为 \(a_{3}=2\)，所以 \(q\neq0\)，并且 \(a_{2}=\frac{2}{q}\)，\(a_{4}=2q\). 代入已知条件得
\(\frac{2}{q}+2q=\frac{20}{3}\)，即 \(3q^{2}-10q+3=0\). 解得 \(q=3\) 或 \(q=\frac{1}{3}\).

当 \(q=3\) 时，\(a_{n}=a_{3}q^{n-3}=2\cdot3^{n-3}\)；当 \(q=\frac{1}{3}\) 时，\(a_{n}=a_{3}q^{n-3}=2\left(\frac{1}{3}\right)^{n-3}=18\left(\frac{1}{3}\right)^{n-1}\). 所以通项公式为 \(a_n=2\cdot3^{n-3}\) 或 \(a_n=18\left(\frac{1}{3}\right)^{n-1}\).
\end{solution}

\begin{problem}
\(\triangle ABC\) 的三个内角为 \(A\)，\(B\)，\(C\)，求当 \(A\) 为何值时，\(\cos A + 2\cos \frac{B + C}{2}\) 取得最大值，并求出这个最大值.
\end{problem}

\begin{solution}
由三角形内角和 \(A+B+C=\pi\)，得 \(\frac{B+C}{2}=\frac{\pi}{2}-\frac{A}{2}\)，从而 \(\cos\frac{B+C}{2}=\sin\frac{A}{2}\). 令 \(t=\sin\frac{A}{2}\)，则 \(0<t<1\)，所求式为
\(\cos A+2\cos\frac{B+C}{2}=1-2t^{2}+2t=\frac{3}{2}-2\left(t-\frac{1}{2}\right)^{2}\leqslant\frac{3}{2}\).

等号成立当且仅当 \(t=\frac{1}{2}\). 因为 \(0<\frac{A}{2}<\frac{\pi}{2}\)，所以 \(\frac{A}{2}=\frac{\pi}{6}\)，即 \(A=\frac{\pi}{3}\). 最大值为 \(\frac{3}{2}\).
\end{solution}

\begin{problem}
\(A\)，\(B\) 是治疗同一种疾病的两种药，用若干试验组进行对比试验. 每个试验组由 \(4\) 只小白鼠组成，其中 \(2\) 只服用 \(A\)，另 \(2\) 只服用 \(B\)，然后观察疗效. 若在一个试验组中，服用 \(A\) 有效的小白鼠的只数比服用 \(B\) 有效的多，就称该试验组为甲类组. 设每只小白鼠服用 \(A\) 有效的概率为 \(\frac{2}{3}\)，服用 \(B\) 有效的概率为 \(\frac{1}{2}\).

\begin{enumerate}
\item 求一个试验组为甲类组的概率.
\item 观察 \(3\) 个试验组，求这 \(3\) 个试验组中至少有一个甲类组的概率.
\end{enumerate}
\end{problem}

\begin{solution}
\begin{enumerate}
\item 设一个试验组中服用 \(A\) 有效的小白鼠数为 \(X\)，服用 \(B\) 有效的小白鼠数为 \(Y\). 则 \(X\) 与 \(Y\) 相互独立，且 \(X\sim B\left(2,\frac{2}{3}\right)\)，\(Y\sim B\left(2,\frac{1}{2}\right)\). 甲类组的条件是 \(X>Y\). 因此
\(P(Y=0)=\frac{1}{4}\)，\(P(X\geqslant1)=1-\left(\frac{1}{3}\right)^{2}=\frac{8}{9}\)；\(P(Y=1)=\frac{1}{2}\)，\(P(X=2)=\left(\frac{2}{3}\right)^{2}=\frac{4}{9}\). 当 \(Y=2\) 时不可能有 \(X>Y\). 所以一个试验组为甲类组的概率为
\(P(X>Y)=\frac{1}{4}\times\frac{8}{9}+\frac{1}{2}\times\frac{4}{9}=\frac{4}{9}\).

\item 不同试验组相互独立，且每组为甲类组的概率均为 \(\frac{4}{9}\). 因此三组中至少有一个甲类组的概率为
\(1-\left(1-\frac{4}{9}\right)^{3}=1-\left(\frac{5}{9}\right)^{3}=\frac{604}{729}\).
\end{enumerate}
\end{solution}

\begin{problem}
如图，\(l_{1}\)，\(l_{2}\) 是互相垂直的异面直线，\(MN\) 是它们的公垂线段. 点 \(A\)，\(B\) 在 \(l_{1}\) 上，\(C\) 在 \(l_{2}\) 上，\(AM = MB = MN\).

\begin{enumerate}
\item 证明：\(AC \perp NB\).
\item 若 \(\angle ACB = 60^{\circ}\)，求 \(NB\) 与平面 \(ABC\) 所成角的余弦值.
\end{enumerate}

\bitmapfigure[max width=0.46\linewidth]{img/2006/national_paper_1_liberal/q20_fig1.png}
\end{problem}

\begin{solution}
\begin{enumerate}
\item 设 \(AM=MB=MN=d>0\). 建立直角坐标系，使 \(M=(0,0,0)\)，\(l_{1}\) 为 \(x\) 轴，\(N=(0,d,0)\)，且 \(l_{2}\) 为过 \(N\) 且平行于 \(z\) 轴的直线. 于是可设 \(A=(-d,0,0)\)，\(B=(d,0,0)\)，\(C=(0,d,h)\). 有
\(\overrightarrow{AC}=(d,d,h)\)，\(\overrightarrow{NB}=(d,-d,0)\)，所以 \(\overrightarrow{AC}\cdot\overrightarrow{NB}=d^{2}-d^{2}=0\)，从而 \(AC\perp NB\).

\item 由 \(\angle ACB=60^{\circ}\)，有 \(\overrightarrow{CA}=(-d,-d,-h)\)，\(\overrightarrow{CB}=(d,-d,-h)\)，故
\(\cos60^{\circ}=\frac{\overrightarrow{CA}\cdot\overrightarrow{CB}}{|CA|\,|CB|}=\frac{h^{2}}{2d^{2}+h^{2}}\). 解得 \(h^{2}=2d^{2}\).

设 \(NB\) 与平面 \(ABC\) 所成的锐角为 \(\varphi\). 平面 \(ABC\) 的一个法向量为
\(\overrightarrow{AB}\times\overrightarrow{AC}=(2d,0,0)\times(d,d,h)=(0,-2dh,2d^{2})\). 因而
\(\sin\varphi=\frac{|\overrightarrow{NB}\cdot(\overrightarrow{AB}\times\overrightarrow{AC})|}{|NB|\,|\overrightarrow{AB}\times\overrightarrow{AC}|}=\frac{|h|}{\sqrt{2}\sqrt{d^{2}+h^{2}}}=\frac{1}{\sqrt{3}}\).
所以 \(\cos\varphi=\sqrt{1-\frac{1}{3}}=\frac{\sqrt{6}}{3}\).
\end{enumerate}
\end{solution}

\begin{problem}
设 \(P\) 是椭圆 \(\frac{x^2}{a^2} + y^2 = 1\)（\(a > 1\)）短轴的一个端点，\(Q\) 为椭圆上一个动点，求 \(|PQ|\) 的最大值.
\end{problem}

\begin{solution}
由于椭圆关于 \(x\) 轴对称，不妨取短轴端点 \(P=(0,1)\)，设 \(Q=(x,y)\). 由椭圆方程，\(x^{2}=a^{2}\left(1-y^{2}\right)\)，且 \(-1\leqslant y\leqslant1\). 因此
\(|PQ|^{2}=x^{2}+(y-1)^{2}=(1-a^{2})y^{2}-2y+a^{2}+1\).

令 \(b=a^{2}-1>0\)，则 \(|PQ|^{2}=-by^{2}-2y+a^{2}+1\)，其关于 \(y\) 的对称轴为 \(y_{0}=-\frac{1}{b}=-\frac{1}{a^{2}-1}\).

当 \(a\geqslant\sqrt{2}\) 时，\(y_{0}\in[-1,0)\)，最大值在 \(y=y_{0}\) 处取得，并且
\(|PQ|^{2}_{\max}=\frac{a^{4}}{a^{2}-1}\)，所以 \(|PQ|_{\max}=\frac{a^{2}}{\sqrt{a^{2}-1}}\).

当 \(1<a<\sqrt{2}\) 时，\(y_{0}<-1\)，上式在 \([-1,1]\) 上递减，最大值在 \(y=-1\) 处取得，此时 \(|PQ|^{2}=4\)，所以 \(|PQ|_{\max}=2\).
\end{solution}

\begin{problem}
设 \(a\) 为实数，函数 \(f(x) = x^{3} - ax^{2} + (a^{2} - 1)x\) 在 \((- \infty, 0)\) 和 \((1, +\infty)\) 都是增函数，求 \(a\) 的取值范围.
\end{problem}

\begin{solution}
\(f'(x)=3x^{2}-2ax+a^{2}-1\). 先考察区间 \((-\infty,0)\). 若 \(a>0\)，则对 \(x<0\) 有 \(f''(x)=6x-2a<0\)，所以 \(f'(x)\) 在该区间递减，且其下确界为 \(f'(0)=a^{2}-1\)，故需 \(a\geqslant1\). 若 \(a<0\)，二次函数 \(f'(x)\) 的顶点 \(x=\frac{a}{3}\) 位于 \((-\infty,0)\) 内，需其最小值非负，即 \(\frac{2a^{2}}{3}-1\geqslant0\)，从而 \(a\leqslant-\frac{\sqrt{6}}{2}\). 当 \(a=0\) 时，取 \(x=-\frac{1}{2}\) 可得 \(f'(x)=-\frac{1}{4}<0\)，不满足增性.

对上述两类参数考察 \((1,+\infty)\). 当 \(a\leqslant-\frac{\sqrt{6}}{2}\) 时，\(f''(x)=6x-2a>0\)，故 \(f'(x)\geqslant f'(1)=(a-1)^{2}+1>0\). 当 \(a\geqslant1\) 且 \(a\leqslant3\) 时，\(f''(x)\geqslant0\)（在 \(x>1\) 上），仍有 \(f'(x)\geqslant f'(1)>0\)；当 \(a>3\) 时，\(f'(x)\) 在 \(x=\frac{a}{3}\) 处取最小值 \(\frac{2a^{2}}{3}-1>0\). 因此这两类参数均满足第二个区间上的增性.

所以 \(a\) 的取值范围为 \(\left(-\infty,-\frac{\sqrt{6}}{2}\right]\cup[1,+\infty)\).
\end{solution}
